%%=============================================================================
%% Samenvatting
%%=============================================================================

% TODO: De "abstract" of samenvatting is een kernachtige (~ 1 blz. voor een
% thesis) synthese van het document.
%
% Een goede abstract biedt een kernachtig antwoord op volgende vragen:
%
% 1. Waarover gaat de bachelorproef?
% 2. Waarom heb je er over geschreven?
% 3. Hoe heb je het onderzoek uitgevoerd?
% 4. Wat waren de resultaten? Wat blijkt uit je onderzoek?
% 5. Wat betekenen je resultaten? Wat is de relevantie voor het werkveld?
%
% Daarom bestaat een abstract uit volgende componenten:
%
% - inleiding + kaderen thema
% - probleemstelling
% - (centrale) onderzoeksvraag
% - onderzoeksdoelstelling
% - methodologie
% - resultaten (beperk tot de belangrijkste, relevant voor de onderzoeksvraag)
% - conclusies, aanbevelingen, beperkingen
%
% LET OP! Een samenvatting is GEEN voorwoord!

%%---------- Nederlandse samenvatting -----------------------------------------
%
% TODO: Als je je bachelorproef in het Engels schrijft, moet je eerst een
% Nederlandse samenvatting invoegen. Haal daarvoor onderstaande code uit
% commentaar.
% Wie zijn bachelorproef in het Nederlands schrijft, kan dit negeren, de inhoud
% wordt niet in het document ingevoegd.

\IfLanguageName{english}{%
\selectlanguage{dutch}
\chapter*{Samenvatting}

Deze bachelorproef onderzoekt hoe een gefaseerde implementatie van een open-source ERP-systeem een servicegerichte Belgische KMO kan ondersteunen bij procesintegratie, vermindering van dubbele gegevensinvoer en voorbereiding op elektronische facturatie via Peppol. De case van Hanuman BV, een onderneming actief in de verkoop, installatie, onderhoud en herstelling van professionele koffieapparatuur, vormt hierbij de praktische context.

De huidige werkwijze binnen Hanuman BV is gebaseerd op een combinatie van Excel, e-mail en afzonderlijke administratieve documenten. Deze gefragmenteerde aanpak leidt tot inefficiënties zoals dubbele gegevensinvoer, beperkte documenttraceerbaarheid, manuele voorraadopvolging en foutgevoelige facturatie. Daarnaast vormt de Belgische verplichting rond gestructureerde B2B-e-facturatie een bijkomende aanleiding om de bestaande processen te evalueren en te verbeteren.

De centrale onderzoeksvraag onderzoekt hoe een stapsgewijze migratie naar Odoo KMO's zoals Hanuman BV kan helpen om doorlooptijd, foutgevoeligheid en dubbele gegevensinvoer te verminderen, en tegelijk de voorbereiding op B2B-e-invoicing via Peppol te ondersteunen. Het doel van deze bachelorproef is het ontwerpen en evalueren van een reproduceerbaar migratiekader voor een servicegerichte KMO-context.

De methodologie is gebaseerd op een ontwerpgerichte casestudy. Eerst werd een AS-IS-analyse uitgevoerd op basis van shadowing, documentanalyse en stakeholderinput. Daarna werden requirements opgesteld, een TO-BE-situatie ontworpen en een gefaseerd migratiekader uitgewerkt. Vervolgens werd dit kader gevalideerd via een Proof of Concept in een educatieve Odoo-testomgeving met representatieve stamdata.

De PoC omvatte zes scenario's: een Sales-to-Cash-flow, een aankoop- en voorraadflow, een repair flow, een gecombineerde onderhouds- en onderdelenfactuur, serienummertracking met preventief onderhoud en een Peppol-readinesscontrole. Binnen deze scenario's werd nagegaan of Odoo de kernprocessen van Hanuman BV functioneel kan ondersteunen via standaardmodules zoals Sales, Purchase, Inventory, Repair, Maintenance, Accounting/Invoicing en Contacts.

De resultaten tonen aan dat Odoo binnen de gecontroleerde testomgeving de onderzochte kernprocessen geïntegreerd kan ondersteunen. Offertes, verkooporders, leveringen en facturen worden beter traceerbaar, voorraadmutaties worden gekoppeld aan aankoop, verkoop en repair, en serienummertracking maakt after-sales opvolging en onderhoudsplanning mogelijk. De KPI-evaluatie wijst indicatief op minder manuele stappen, snellere facturatie en betere documenttraceerbaarheid, maar deze cijfers zijn scenario-gebaseerd en vormen geen statistisch bewezen productieresultaten.

Voor Peppol toont de PoC aan dat Odoo een basis kan bieden voor e-invoicing-readiness. Bedrijfsgegevens, klantgegevens, btw-instellingen, bankgegevens en factuurvelden konden centraal worden beheerd. Daarnaast werd in de testomgeving een Peppol-demo-optie met UBL BIS 3 XML-bijlage zichtbaar. Een effectieve productie-uitwisseling via een gecertificeerd Peppol Access Point viel buiten de scope van dit onderzoek.

Op basis van deze bevindingen kan geconcludeerd worden dat een gefaseerde implementatie van Odoo haalbaar en zinvol is voor een servicegerichte Belgische KMO zoals Hanuman BV. De meerwaarde ligt vooral in procesintegratie, centraal databeheer, betere traceerbaarheid en voorbereiding op digitale facturatie. Tegelijk blijft succes afhankelijk van datakwaliteit, change management, gebruikersacceptatie en verdere validatie in een productieomgeving.

\selectlanguage{english}
}{}

%%---------- Samenvatting -----------------------------------------------------
% De samenvatting in de hoofdtaal van het document

\chapter*{\IfLanguageName{dutch}{Samenvatting}{Abstract}}

Deze bachelorproef onderzoekt hoe een gefaseerde implementatie van een open-source ERP-systeem een servicegerichte Belgische KMO kan ondersteunen bij procesintegratie, vermindering van dubbele gegevensinvoer en voorbereiding op elektronische facturatie via Peppol. De case van Hanuman BV, een onderneming actief in de verkoop, installatie, onderhoud en herstelling van professionele koffieapparatuur, vormt hierbij de praktische context.

De huidige werkwijze binnen Hanuman BV is gebaseerd op een combinatie van Excel, e-mail en afzonderlijke administratieve documenten. Deze gefragmenteerde aanpak leidt tot inefficiënties zoals dubbele gegevensinvoer, beperkte documenttraceerbaarheid, manuele voorraadopvolging en foutgevoelige facturatie. Daarnaast vormt de Belgische verplichting rond gestructureerde B2B-e-facturatie een bijkomende aanleiding om de bestaande processen te evalueren en te verbeteren.

De centrale onderzoeksvraag onderzoekt hoe een stapsgewijze migratie naar Odoo KMO's zoals Hanuman BV kan helpen om doorlooptijd, foutgevoeligheid en dubbele gegevensinvoer te verminderen, en tegelijk de voorbereiding op B2B-e-invoicing via Peppol te ondersteunen. Het doel van deze bachelorproef is het ontwerpen en evalueren van een reproduceerbaar migratiekader voor een servicegerichte KMO-context.

De methodologie is gebaseerd op een ontwerpgerichte casestudy. Eerst werd een AS-IS-analyse uitgevoerd op basis van shadowing, documentanalyse en stakeholderinput. Daarna werden requirements opgesteld, een TO-BE-situatie ontworpen en een gefaseerd migratiekader uitgewerkt. Vervolgens werd dit kader gevalideerd via een Proof of Concept in een educatieve Odoo-testomgeving met representatieve stamdata.

De PoC omvatte zes scenario's: een Sales-to-Cash-flow, een aankoop- en voorraadflow, een repair flow, een gecombineerde onderhouds- en onderdelenfactuur, serienummertracking met preventief onderhoud en een Peppol-readinesscontrole. Binnen deze scenario's werd nagegaan of Odoo de kernprocessen van Hanuman BV functioneel kan ondersteunen via standaardmodules zoals Sales, Purchase, Inventory, Repair, Maintenance, Accounting/Invoicing en Contacts.

De resultaten tonen aan dat Odoo binnen de gecontroleerde testomgeving de onderzochte kernprocessen geïntegreerd kan ondersteunen. Offertes, verkooporders, leveringen en facturen worden beter traceerbaar, voorraadmutaties worden gekoppeld aan aankoop, verkoop en repair, en serienummertracking maakt after-sales opvolging en onderhoudsplanning mogelijk. De KPI-evaluatie wijst indicatief op minder manuele stappen, snellere facturatie en betere documenttraceerbaarheid, maar deze cijfers zijn scenario-gebaseerd en vormen geen statistisch bewezen productieresultaten.

Voor Peppol toont de PoC aan dat Odoo een basis kan bieden voor e-invoicing-readiness. Bedrijfsgegevens, klantgegevens, btw-instellingen, bankgegevens en factuurvelden konden centraal worden beheerd. Daarnaast werd in de testomgeving een Peppol-demo-optie met UBL BIS 3 XML-bijlage zichtbaar. Een effectieve productie-uitwisseling via een gecertificeerd Peppol Access Point viel buiten de scope van dit onderzoek.

Op basis van deze bevindingen kan geconcludeerd worden dat een gefaseerde implementatie van Odoo haalbaar en zinvol is voor een servicegerichte Belgische KMO zoals Hanuman BV. De meerwaarde ligt vooral in procesintegratie, centraal databeheer, betere traceerbaarheid en voorbereiding op digitale facturatie. Tegelijk blijft succes afhankelijk van datakwaliteit, change management, gebruikersacceptatie en verdere validatie in een productieomgeving.